\setcounter{section}{53}
\section{Equipos personales}

Nombre completo: Dispositivos personales de PC y dispositivos móviles. Conectividad de dispositivos personales. Medidas de seguridad y gestión para equipos personels y dispositivos móviles.

\localtableofcontents

\subsection{Dispositivos personales PC}
Mira macho, paso de contarte esto. Así que te interese las NPU>

\subsubsection{Neural Processing Unit}
Aceleradores diseñados específicamente para ejecutar algoritmos de inteligencia artificial. Estámn optimizados para inferencia mas que entrenamiento. 

\subsubsection{Tensor Processing Unit}
Los Tensores son estructuras de datos para el aprendizaje profundo. Están orientadas para entornos clod y CPD. Se usan sobre todo para el entrenamiento masivo.

\subsection{Dispositivos móviles}

\textit{Mira macho}

\subsection{Conectividad de dispositivos personales}

\subsubsection{Wifi direct}
como el bluetooth, pero por wifi.

\subsection{Medidas de seguridad}
Las herramientas de seguridad y gestión se han convertido en un elemento esencial.

\begin{description}
    \item[MDM] Mobile Device Management, centrado en la gestión de dispositivos.
    \item[EMM] Enterprise Mobile Management, pone el foco en la seguridad de la información corporativa y aplicaciones. En lugar de proteger el dispositivo en sí, se reconocen políticas de BYOD protegiendo la parte empresarial.
    \item[UEM] Unified Endpoint Management incorpora todos los tipos de dispositivo en una única herramienta. Permite múltiples sistemas y dispositivos.
\end{description}

Algiunas funciones a tener en cuenta son:
\begin{itemize}
    \item Inventario
    \item Monitorización
    \item Cifrado
    \item Restricciones HW y SW
    \item Actualización de SW y mantenimiento remoto
    \item Borrado remoto
    \item Localización
    \item Mobile Threat Prevention & Mobile Threat Defense
    \item Copias de seguridad
\end{itemize}

Se tienen en cuenta los siguentes modelos de despliegue:
\begin{description}
    \item[COSU Corporate Owned Single uUse] limita un dispositivo a una sola aplicación dedicada.
    \item[COBO Corporate Owned, Business Only] Limitan el uso empresarial, prohibiendo o restringiendo otros usos. Pueden tener aplicaciones en whitelist y límites de configuración.
    \item[COPE Corporate Owned, Personally Enabled] Permite a los empleados el uso personal del dispositivo. Permite igualmente ciertas restricciones en el uso, aplicaciones, conexiones, etc.
    \item[CYOU Choose Your Own DEvice] proporciona al usuario una selección de dispositivos aprobados. Permitiendo adaptación para el usuario manteniendo medidas de seguridad.
    \item[BYOD] Permite al empleado utilizar su propio dispositivo personal. Reduce el coste de proveer de dispositivos empleados. También supone un riesgo mayor al no poder forzar todas las políticas de seguridad. Igualmente si se pueden forzar temas de antivirus, contraseñas, registro obligatorio, acceso condicional, limpieza remota o sanciones en caso de incumplimiento.
\end{description}

La mobilidad tiene ciertos riesgos. Existiendo gúias del Centro Criptográfico Nacional:
\begin{itemize}
    \item CCN-STIC 827 para BYOD
    \item CCN-STIC 450 para evaluar riesgos, amenazas y vulnerabilidades.
    \begin{itemize}
        \item CCN-STIC 453 Para seguridad en Android
        \item CCN-STIC 454 Para seguridad en iPad
        \item CCN-STIC 455 Para seguridad en iPhone
        \item CCN-STIC 457 Para seguridad en  dispositivos móviles MDM
    \end{itemize}
\end{itemize}